\subsection{Occupancy collapse of the attachment polynomial}
\label{subsection-attachment-occupancy-collapse}

Fix labelled canonical atoms $A_1,\ldots,A_k$, retain all labels on their
vertices, and put
\[
w_i=|V(A_i)|,
\qquad
P=\prod_iw_i,
\qquad
v=\sum_iw_i-k+1.
\]
Thus $v$ is the order of every connected quotient assembly.  Let $N_t$ count
capacity-respecting connected attachment quotients with exactly $t$ shared
points.  The preceding enumeration gives
\[
N_t=P\,S(k-1,t)(v-1)_{t-1}.
\]

\begin{theorem}[Attachment occupancy collapse]
\label{theorem-attachment-occupancy-collapse}
For $k\ge2$,
\[
N_t=\frac Pv(v)_tS(k-1,t),
\qquad
\sum_{t=1}^{k-1}N_t=Pv^{k-2}.
\tag{A.1}
\]
Consequently, under the uniform distribution on attachment quotients, the
number $T$ of shared points satisfies
\[
\Pr(T=t)=\frac{(v)_tS(k-1,t)}{v^{k-1}}.
\tag{A.2}
\]
This is the occupancy law for $k-1$ labelled balls placed uniformly into $v$
labelled boxes.
\end{theorem}
\begin{proof}
Use $(v-1)_{t-1}=(v)_t/v$ and
$x^n=\sum_tS(n,t)(x)_t$.  Equivalently, normalize the bipartite Prüfer code by
choosing one anchor vertex in each atom.  The remaining code alphabet has
$1+\sum_i(w_i-1)=v$ symbols, and one code position may be fixed, leaving
$v^{k-2}$ words.
\end{proof}

For an excess profile
$\lambda=(\lambda_1,\ldots,\lambda_t)\vdash k-1$, let $m_r$ count its parts
equal to $r$.  The full profile law is
\[
\Pr(\lambda(F)=\lambda)
=
\frac{(v)_t(k-1)!}
 {v^{k-1}\prod_j\lambda_j!\prod_rm_r!}.
\tag{A.3}
\]
Thus the complete shared-point multiplicity profile is the load profile in the
same occupancy model.  If $C_r$ counts shared points of multiplicity $r+1$,
then
\[
\mathbb EC_r
=v\binom{k-1}{r}v^{-r}
 \left(1-\frac1v\right)^{k-1-r}.
\tag{A.4}
\]

Put
\[
a=\left(1-\frac1v\right)^{k-1},
\qquad
b=\left(1-\frac2v\right)^{k-1}.
\]
Then
\[
\mathbb ET=v(1-a),
\qquad
\operatorname{Var}(T)=va+v(v-1)b-v^2a^2.
\tag{A.5}
\]
More generally,
\[
\mathbb E(T)_j
=(v)_j\sum_{\ell=0}^j(-1)^\ell\binom j\ell
\left(1-\frac\ell v\right)^{k-1}.
\tag{A.6}
\]

\begin{theorem}[Real-rooted attachment polynomial]
\label{theorem-real-rooted-attachment-polynomial}
The polynomial
\[
\Phi_{\mathbf A}(z)=\sum_{t=1}^{k-1}N_tz^t
\]
has one simple zero at $0$ and all its other zeros are simple and negative.
The positive coefficients are strictly ultra log-concave.  Moreover, there
exist independent Bernoulli variables $B_1,\ldots,B_{k-2}$ such that
\[
T\stackrel d=1+\sum_{j=1}^{k-2}B_j.
\tag{A.7}
\]
\end{theorem}
\begin{proof}
Set
\[
\Omega_{h,v}(z)=\sum_t(v)_tS(h,t)z^t.
\]
The Stirling recurrence gives
\[
\Omega_{h+1,v}=vz\Omega_{h,v}+z(1-z)\Omega'_{h,v}.
\tag{A.8}
\]
Starting from $\Omega_{1,v}=vz$, evaluation at the old negative zeros shows
alternating signs, while the signs just left of zero and at negative infinity
supply the two endpoint intervals.  Induction gives simple negative zeros and
strict interlacing.  Newton's inequalities give strict ultra log-concavity.
Factoring the normalized probability generating polynomial over its negative
zeros gives (A.7).
\end{proof}

The occupancy interpretation distinguishes two asymptotic regimes.  For fixed
$k$ and $v\to\infty$, binary attachments dominate and
\[
\Pr(T=k-1)=\frac{(v)_{k-1}}{v^{k-1}}
=1-\frac{\binom{k-1}{2}}v+O_k(v^{-2}).
\tag{A.9}
\]
If instead $h=k-1\to\infty$ and $v/h\to\gamma\ge1$, then
\[
\frac{\mathbb ET}{h}\to\gamma(1-e^{-1/\gamma}),
\qquad
\frac{\operatorname{Var}(T)}h
\to\gamma e^{-1/\gamma}-(\gamma+1)e^{-2/\gamma},
\tag{A.10}
\]
and the centered, variance-normalized $T$ is asymptotically Gaussian.  For
fixed $r$,
\[
\frac{C_r}{h}
\xrightarrow{\Pr}
\frac{e^{-1/\gamma}}{r!\gamma^{r-1}}.
\tag{A.11}
\]
Hence binary dominance for fixed $k$ and large atom orders does not persist
when the number of fixed-size atoms tends to infinity.
